% Part: lambda-calculus
% Chapter: introduction
% Section: computable-lambda

\documentclass[../../../include/open-logic-section]{subfiles}

\begin{document}

\olfileid{lam}{int}{lrp}
\olsection{Fungsi kang Bisa Diitung iku \usetoken{S}{lambda definable}}

\begin{thm}
\ollabel{thm:computable-lambda}
Saben fungsi parsial kang bisa diitung iku
!!{lambda definable}.
\end{thm}

\begin{proof}
Kita kudu nuduhake yen saben fungsi parsial kang
bisa diitung~$f$ !!{lambda defined} dening sawijining
suku lambda~$F$. Miturut teorema wujud normal
Kleene, cukup nuduhake yen saben fungsi rekursif
primitif !!{lambda defined} dening sawijining suku
lambda, banjur yen kelas fungsi kang
!!{lambda definable} katutup tumrap komposisi
kang trep lan panelusuran tanpa wates. Kanggo
nuduhake yen saben fungsi rekursif primitif
!!{lambda defined} dening sawijining suku lambda,
cukup nuduhake yen fungsi-fungsi wiwitan iku
!!{lambda definable}, lan yen kelas fungsi parsial
kang !!{lambda definable} katutup tumrap komposisi,
rekursi primitif, lan panelusuran tanpa wates.
\end{proof}

Kita bakal nggunakake notasi kang luwih lumrah
supaya bukti sabanjure luwih gampang diwaca.
Umpamane, kita bakal nulis $M(x, y, z)$
tinimbang $Mxyz$. Senajan notasi iki menehi
gambaran tartamtu, elinga yen suku ing
kalkulus lambda tanpa tipe ora nduweni
aritas kang tetep; mula kanggo suku~$M$
kang padha, nulis $M(x,y)$ utawa
$M(x,y,z,w)$ padha-padha lumrah.
Nanging notasi iki nuduhake yen kita lagi
nganggep $M$ minangka fungsi telung variabel,
lan mbantu njlentrehake maksud definisi.
Kanthi cara kang padha, kita bakal nyebut
``netepake $M$ kanthi $M(x,y,z) = \dots$''
tinimbang ``netepake $M$ kanthi $M =
\lambd[x][\lambd[y][\lambd[z][\dots]]]$.''

\end{document}
